\documentclass{stex}

\libusepackage{naproche}
\libusepackage{copyright}
\libinput{preamble}

\begin{document}
\begin{smodule}{first-order-semantics}

\importmodule[libraries/meta]{preliminaries.ftl}

\begin{sfragment}[id=sec:first-order-semantics]{First-Order Semantics}
   \symdecl*{formula}
  \symdef{folsem}[name=first-order semantics,args=1,op=[\![\cdot]\!]]{\comp{[\![}#1\comp{]\!]}}
  %
  \begin{sparagraph}[style=symdoc,for={first-order semantics,folsem}]
    \vardef{Evar}{E}
    %
    \Naproche translates \ForTheL expressions to (a variant of)
    \sr{formula}{first-order formulas} which we refer to as the
    \definame{first-order semantics} of \ForTheL.
    We denote the first-order translation of an expression $\Evar$
    by $\defnotation\folsem\Evar$.
  \end{sparagraph}

  \begin{sfragment}[id=sec:tags]{Tags}
    \symdecl*{tag}
    \symdef{DigTag}{\langle\textsf{dig}\rangle}
    \symdef{DigMultiSubjectTag}{\langle\textsf{dig multi-subject}\rangle}
    \symdef{DigMultiPairwiseTag}{\langle\textsf{dig multi-pairwise}\rangle}
    \symdef{HeadTermTag}{\langle\textsf{head term}\rangle}
    \symdef{InductionHypothesisTag}{\langle\textsf{induction hypothesis}\rangle}
    \symdef{CaseHypothesisTag}{\langle\textsf{case hypothesis}\rangle}
    \symdef{EqualityChainTag}{\langle\textsf{equality chain}\rangle}
    \symdef{GenericMarkTag}{\langle\textsf{generic mark}\rangle}
    \symdef{EvaluationTag}{\langle\textsf{evaluation}\rangle}
    \symdef{ConditionTag}{\langle\textsf{condition}\rangle}
    \symdef{DefinedTag}{\langle\textsf{defined}\rangle}
    \symdef{DomainTag}{\langle\textsf{domain}\rangle}
    \symdef{ReplacementTag}{\langle\textsf{replacement}\rangle}
    \symdef{DomainTaskTag}{\langle\textsf{domain task}\rangle}
    \symdef{ExistenceTaskTag}{\langle\textsf{existence task}\rangle}
    \symdef{UniquenessTaskTag}{\langle\textsf{uniqueness task}\rangle}
    \symdef{ChoiceTaskTag}{\langle\textsf{choice task}\rangle}
    %
    \begin{sparagraph}[style=symdoc,for={tag,DigTag,DigMultiSubjectTag,DigMultiPairwiseTag,HeadTermTag,InductionHypothesisTag,
      CaseHypothesisTag,EqualityChainTag,GenericMarkTag,EvaluationTag,ConditionTag,DefinedTag,DomainTag,ReplacementTag,
      DomainTaskTag,ExistenceTaskTag,UniquenessTaskTag,ChoiceTaskTag}]
      \Sns{formula} can be annotated with one of the following \definame[post=s]{tag}:
      \begin{enumerate}
        \item $\defnotation\DigTag$
        \item $\defnotation\DigMultiSubjectTag$
        \item $\defnotation\DigMultiPairwiseTag$
        \item $\defnotation\HeadTermTag$
        \item $\defnotation\InductionHypothesisTag$
        \item $\defnotation\CaseHypothesisTag$
        \item $\defnotation\EqualityChainTag$
        \item $\defnotation\GenericMarkTag$
        \item $\defnotation\EvaluationTag$
        \item $\defnotation\ConditionTag$
        \item $\defnotation\DefinedTag$
        \item $\defnotation\DomainTag$
        \item $\defnotation\ReplacementTag$
        \item $\defnotation\DomainTaskTag$
        \item $\defnotation\ExistenceTaskTag$
        \item $\defnotation\UniquenessTaskTag$
        \item $\defnotation\ChoiceTaskTag$
      \end{enumerate}
    \end{sparagraph}

    $\GenericMarkTag$, $\EvaluationTag$, $\ConditionTag$, $\DefinedTag$, $\DomainTag$ and $\ReplacementTag$ \sns{tag}
    are used to mark certain parts of map definitions.
    $\DomainTaskTag$, $\ExistenceTaskTag$, $\UniquenessTaskTag$ and $\ChoiceTaskTag$ \sns{tag} are used to mark
    certain parts of well-definedness proof tasks for maps.
  \end{sfragment}

  \begin{sfragment}[id=sec:formulas]{Formulas}
    \symdef{foltrue}{\top}
    \symdef{folfalse}{\bot}
    \symdef{folthis}{\textsf{this}}
    \symdef{folneg}[args=1]{\mathop{\maincomp\neg}#1}
    \symdef{folconj}[args=a]{\argsep{#1}{\mathbin{\maincomp\wedge}}}
    \symdef{foldisj}[args=a]{\argsep{#1}{\mathbin{\maincomp\vee}}}
    \symdef{folimpl}[args=2]{#1\mathbin{\maincomp\Rightarrow}#2}
    \symdef{folequiv}[args=a]{\argsep{#1}{\mathbin{\maincomp\Leftrightarrow}}}
    \symdef{folforall}[args=1]{\maincomp\forall#1}
    \symdef{folexists}[args=1]{\maincomp\exists#1}
    \symdef{foltag}[args=2]{#1\mathrel{\maincomp{::}}#2}
    \symdef{foldbindex}[args=1,name=De Brujin index]{#1}
    %
    \begin{sparagraph}[style=symdoc,for={formula,foltrue,folfalse,folneg,folconj,foldisj,folimpl,folequiv,folforall,folexists,foltag,foldbindex,folsem}]
      \vardef{phivar}{\varphi}
      \vardef{psivar}{\psi}
      \vardef{Tvar}{T}
      \vardef{ivar}{i}
      \vardef{xvar}{x}
      \vardef{tvar}[args=a]{\maincomp t\dobrackets{#1}}
      \vardef{nvar}{n}
      \varseq{phiseq}{1,\ellipses,\nvar}{\maincomp\varphi_{#1}}
      %
      \Definame[post=s]{formula}, the target of the \sn{first-order semantics} of \Naproche,
      are recursively defined as follows:
      \begin{enumerate}
        \item $\defnotation\foltrue$ is a \sn{formula};
        \item $\defnotation\folfalse$ is a \sn{formula};
        \item $\defnotation\folthis$ is a \sn{formula};
        \item if $\phivar$ is a \sn{formula},
          then $\defnotation\folneg\phivar$ is a \sn{formula};
        \item if $\phivar$ and $\psivar$ are \sns{formula},
          then $\defnotation\folconj{\phivar,\psivar}$ is a \sn{formula};
        \item if $\phivar$ and $\psivar$ are \sns{formula},
          then $\defnotation\foldisj{\phivar,\psivar}$ is a \sn{formula};
        \item if $\phivar$ and $\psivar$ are \sns{formula},
          then $\defnotation\folimpl{\phivar}{\psivar}$ is a \sn{formula};
        \item if $\phivar$ and $\psivar$ are \sns{formula},
          then $\defnotation\folequiv{\phivar,\psivar}$ is a \sn{formula};
        \item if $\phivar$ is a \sn{formula},
          then $\defnotation\folforall\phivar$ is a \sn{formula};
        \item if $\phivar$ is a \sn{formula},
          then $\defnotation\folexists\phivar$ is a \sn{formula};
        \item if $\Tvar$ is a tag and $\phivar$ a \sn{formula},
          then $\defnotation\foltag\Tvar\phivar$ is a \sn{formula};
        \item if $\ivar\in\mathbb N$,
          then $\defnotation\foldbindex\ivar$ is a \sn{formula},
          commonly called a \definame{De Brujin index};
        \item if $\xvar$ is a variable name,
          then $\xvar$ is a \sn{formula};
        \item if $\tvar!$ is a term name and $\phiseq!$ are \sns{formula},
          then $\tvar\phiseq$ is a \sn{formula}.
      \end{enumerate}
    \end{sparagraph}
  \end{sfragment}
\end{sfragment}

\ifinputref\else\printlicense[CcByNcSa]{Marcel Schütz (2026)}\fi

\end{smodule}
\end{document}
